% Original worked bridge: Bloch transform, eigenvalues and Berry connection.
\subsection*{Worked derivation: from SSH hopping to a winding number}
We now unpack the short step from the real-space hopping
Hamiltonian to its two-band Bloch matrix. Choose a ring of $N$
unit cells, momentum values $k=2\pi m/N$, and the Fourier
conventions of Eq.~\eqref{eq:ssh-real}:
\begin{equation}
\op a_n=\frac1{\sqrt N}\sum_k e^{\ii kn}\op a_k,\qquad
\op b_n=\frac1{\sqrt N}\sum_q e^{\ii qn}\op b_q.
\end{equation}
To transform the intracell term, substitute both expansions:
\begin{align}
\sum_n \op a_n^\dagger\op b_n
&=\frac1N\sum_{n,k,q}
 e^{-\ii kn}e^{\ii qn}\op a_k^\dagger\op b_q\notag\\
&=\sum_{k,q}\left[\frac1N\sum_n
 e^{\ii(q-k)n}\right]\op a_k^\dagger\op b_q
=\sum_k\op a_k^\dagger\op b_k,
\end{align}
where the term in square brackets is $\delta_{kq}$.
For the intercell term a shift of the position index matters:
\begin{align}
\sum_n\op a_{n+1}^\dagger\op b_n
&=\frac1N\sum_{n,k,q}
 e^{-\ii k(n+1)}e^{\ii qn}
 \op a_k^\dagger\op b_q\notag\\
&=\sum_k e^{-\ii k}\op a_k^\dagger\op b_k.
\end{align}
Together with the Hermitian conjugate,
\begin{equation}
\op H=\sum_k\begin{pmatrix}
\op a_k^\dagger&\op b_k^\dagger
\end{pmatrix}
\underbrace{\begin{pmatrix}
0&t_1+t_2e^{-\ii k}\\
t_1+t_2e^{\ii k}&0
\end{pmatrix}}_{h(k)}
\begin{pmatrix}\op a_k\\\op b_k\end{pmatrix}.
\end{equation}
There is now a precise reason for the exponent sign in
the upper-right matrix entry.

\paragraph{Derive the two energies and locate the gap.}
Let $q=t_1+t_2e^{-\ii k}$. The characteristic polynomial is
\begin{equation}
\det[h(k)-EI]
=\begin{vmatrix}-E&q\\q^*&-E\end{vmatrix}
=E^2-|q|^2.
\end{equation}
Compute the modulus explicitly:
\begin{equation}
\begin{aligned}
|q|^2
&=(t_1+t_2e^{-\ii k})(t_1+t_2e^{\ii k})\\
&=t_1^2+t_2^2+t_1t_2(e^{\ii k}+e^{-\ii k})\\
&=t_1^2+t_2^2+2t_1t_2\cos k.
\end{aligned}
\end{equation}
Thus $E_\pm=\pm|q|$. For $t_1,t_2>0$,
the minimum $|q|$ occurs at $k=\pi$, because
$\cos k\geq-1$. Hence
$|q(\pi)|=|t_1-t_2|$ and
the direct spectral gap is $2|t_1-t_2|$.

\paragraph{One explicit lower-band eigenvector.}
Write the \emph{upper-right} off-diagonal element as
$q=|q|e^{-\ii\varphi}$, so that
$q^*=|q|e^{\ii\varphi}$. With
$\varphi=\arg(t_1+t_2e^{\ii k})$, choose
\begin{equation}
\ket{u_-(k)}=\frac1{\sqrt2}
\begin{pmatrix}-e^{-\ii\varphi(k)}\\1\end{pmatrix}.
\end{equation}
Multiplication verifies its lower-band eigenvalue:
\begin{equation}
h(k)\ket{u_-}
=\frac{|q|}{\sqrt2}
\begin{pmatrix}
e^{-\ii\varphi}\\-1
\end{pmatrix}
=-|q|\ket{u_-}.
\end{equation}
Now differentiate the state:
\begin{equation}
\partial_k\ket{u_-}
=\frac1{\sqrt2}
\begin{pmatrix}\ii e^{-\ii\varphi}\partial_k\varphi\\0
\end{pmatrix}.
\end{equation}
Its Berry connection is
\begin{equation}
\begin{aligned}
\mathcal A(k)
&=\ii\braket{u_-(k)}{\partial_k u_-(k)}\\
&=\frac{\ii}{2}
 \bigl(-e^{\ii\varphi}\bigr)
 \bigl(\ii e^{-\ii\varphi}\partial_k\varphi\bigr)
=\frac12\partial_k\varphi.
\end{aligned}
\end{equation}
The map
$(t_1+t_2\cos k,t_2\sin k)$ is a circle of
radius $t_2$ centered at $(t_1,0)$ in the complex plane.
The argument $\varphi$ winds by $2\pi$ if the circle
encloses the origin (when $t_2>t_1$), and by zero
if it does not (when $t_1>t_2$).
Therefore
\begin{equation}
\gamma_{\mathrm{Zak}}
=\oint\mathcal A(k)\dd k
=\frac12[\varphi(\pi)-\varphi(-\pi)]
=\pi\nu\pmod{2\pi}.
\end{equation}

\paragraph{Physical interpretation and limits.}
A nonzero winding requires a spectral gap and a fixed
choice of sublattice/unit-cell convention. At $t_1=t_2$,
$q(\pi)=0$, so the phase becomes undefined exactly where
the gap closes. Adding a staggered onsite $\sigma_z$
term changes the symmetry class, allowing the planar loop
to move away from the origin in three dimensions.
The interactive investigations introduced alongside this derivation allow the bulk winding, Berry phase, and finite-chain spectrum to be checked independently.

\interactiveguide{IC-08-04}{Zak phase and Berry gauge}{After computing the lower-band eigenvector and Berry connection explicitly, change the unit-cell or phase convention. Compare discrete overlap phases with the analytic Zak phase and identify which quantities remain gauge invariant.}
